\section*{第 \ref{ch:b2:open-systems} 章——热力学衡算与开口系统}

\begin{solution}{exo:b2:open-systems:1}
(a) $\qty{2e-4}{m^3/s}/\qty{7.85e-5}{m^2} = \qty{2.5}{m/s}$. (b) $(0.1 \times 60 + 0.2
\times 10)/0.3 = \qty{26.7}{\celsius}$. (c) $\dot S_{\text{c}} = 0.1c\ln(299.8/333) +
0.2c\ln(299.8/283) = \qty{4.3}{W/K}$：热量从热水传向冷水。
\end{solution}

\begin{solution}{exo:b2:open-systems:2}
$\sqrt{2 \times 200 \times 10^3 + 50^2} = \qty{634}{m/s}$; $c^2/2c_p = \qty{45}{K}$;
\qty{31}{K}.
\end{solution}

\begin{solution}{exo:b2:open-systems:3}
(a) $T$ 不变，$s$ 增加 $r\ln(P_1/P_2)$。(b) $h = \qty{2576}{kJ/kg}$：
$x = (2576 - 417)/2258 = 0.956$——干度由
出口处过热（或恰为干饱和）蒸气的温度读出。(c) $(256 -
200)/199 = 0.28$。
\end{solution}

\begin{solution}{exo:b2:open-systems:4}
(a) \qty{41.8}{kW}; $41.8 \times 10^3/(1005 \times 20) = \qty{2.1}{kg/s}$, \qty{1.7}{m^3/s}.
(b) \qty{10}{MW}. (c) $v\Delta P = \qty{5.9}{kJ/kg}$; \qty{3}{MW}.
\end{solution}

\begin{solution}{exo:b2:open-systems:5}
(a) 由\cref{thm:b2:open-systems:firstlaw}可知。(b) 流体被
相邻流体推入和推出：该压强功（单位质量为 $Pv$）
与 $u$ 合并计入。(c) \qty{50}{J/kg} 和 \qty{45}{kJ/kg}，相较于 \qty{1}{kJ/kg}。
(d) \qty{981}{J/kg} 与 \qty{4180}{J/kg} 相比：\qty{0.23}{K}。
\end{solution}

\begin{solution}{exo:b2:open-systems:6}
(a) \qty{531}{K}, \qty{239}{kJ/kg}. (b) \qty{299}{kJ/kg}, \qty{590}{K}, $c_p\ln(590/531)
= \qty{106}{J/kg/K}$。(c) 每级在 \qty{393}{K} 处结束：$2 \times 1005 \times 100 =
\qty{201}{kJ/kg}$，节省 $16\%$。(d) $rT\ln 8 = \qty{175}{kJ/kg}$；若
在多级之间进行级间冷却，则过程线紧贴等温线。
\end{solution}

\begin{solution}{exo:b2:open-systems:7}
(a) $x = (6.88 - 0.476)/7.918 = 0.81$, $h_{4s} = \qty{2098}{kJ/kg}$, $w_{\text{s}} =
\qty{1324}{kJ/kg}$. (b) \qty{1125}{kJ/kg}; $h_4 = 2297$, $x_4 = 0.89$. (c) $s_4 =
\qty{7.53}{kJ/kg/K}$：$s_{\text{c}} = \qty{0.65}{kJ/kg/K}$；损失 $T_0s_{\text{c}} =
\qty{200}{kJ/kg}$。(d) 速度为数百米每秒的液滴侵蚀
末级叶片。
\end{solution}

\begin{solution}{exo:b2:open-systems:8}
(a) $h_2 = 399 + 27/0.8 = \qty{433}{kJ/kg}$；$w = \qty{34}{kJ/kg}$。(b) $q_{\text{蒸发}}
= 143$，$q_{\text{冷凝}} = \qty{177}{kJ/kg}$；$\mathrm{COP}_{\text{c}} = 4.2$，
$\mathrm{COP}_{\text{h}} = 5.2$。(c) $6.8$ 和 $7.8$；三分之二。(d) \qty{0.045}{kg/s}，
\qty{1.5}{kW}；蒸发器温度越低，意味着压强越低、
压力比越大（每千克做功更多），且蒸气更轻（
每冲程质量更少）。
\end{solution}

\begin{solution}{exo:b2:open-systems:9}
(a) \qty{30}{\celsius}. (b) $4180[\ln(323/363) + 2\ln(303/283)] = \qty{83}{W/K}$.
(c) \qty{24}{kW}，占 \qty{167}{kW}。(d) 并流时温差
更大，且热流体出口温度不能低于冷流体出口温度：
相同负荷下产生的熵更多。
\end{solution}

\begin{solution}{exo:b2:open-systems:10}
(a) $\dd(mu)/\dd t = \dot m h_0$ 积分得 $mu = mh_0$：$c_vT = c_pT_0$，
$T = \gamma T_0$。(b) \qty{410}{K}。(c) 终态气体是原有气体（温度为 $T_0$）
与被加热至 $\gamma T_0$ 的流入气体的混合物，因而
介于二者之间。(d) 流动功使气体升温；缓慢充入
水浴可将热量导出，使 \qty{20}{\celsius} 时的压强
恰为所需。
\end{solution}

\begin{solution}{exo:b2:open-systems:11}
(a) $x = (7.76 - 0.476)/7.918 = 0.92$, $h = \qty{2367}{kJ/kg}$. (b) $w = 572 + 1111
- 6 = \qty{1677}{kJ/kg}$, $q = 3278 + 628 = \qty{3906}{kJ/kg}$, $\eta = 0.43$
与 $0.40$ 相比。(c) 再热在高温下加入热量，并
提高了吸热的平均温度。(d) 排气更干燥
（$0.92$，而非 $0.81$）。
\end{solution}

\begin{solution}{exo:b2:open-systems:12}
(a) \qty{293}{K}, \qty{596}{K}, \qty{1500}{K}, \qty{737}{K}. (b) $w_{\text{c}} = 305$,
$w_{\text{t}} = 767$, $q = \qty{909}{kJ/kg}$, $\eta = 0.51 = 1 - 12^{-0.286}$. (c)
排气热量高于 \qty{373}{K}：\qty{366}{kJ/kg}；$0.7 \times 0.35 \times 366 =
\qty{90}{kJ/kg}$：$\eta = (462 + 90)/909 = 0.61$。(d) 它在
\qty{1500}{K} 处吸热，并在 \qty{306}{K} 处放热：
这是任何工质所允许的最大温度跨度。
\end{solution}

\begin{solution}{pb:b2:open-systems:1}
\textbf{1.} 压缩机（绝热，$w = \Delta h$），冷凝器（等压，
放热），节流阀（等焓），蒸发器（等压，吸热）。

\textbf{2.} $h_4 = \qty{256}{kJ/kg}$, $x_4 = 0.28$.

\textbf{3.} \qty{143}{kJ/kg}.

\textbf{4.} \qty{27}{} 与 \qty{34}{kJ/kg}；$h_2 = 433$；$T_2 \approx 45 + 7 \approx
\qty{52}{\celsius}$。

\textbf{5.} \qty{177}{kJ/kg}; $\mathrm{COP}_{\text{h}} = 5.2$, $\mathrm{COP}_{\text{c}} = 4.2$;
$5.2 = 4.2 + 1$.

\textbf{6.} $7.8$ 与 $6.8$；三分之二。

\textbf{7.} \qty{0.045}{kg/s}；\qty{1.5}{kW}，在插头处为 \qty{1.7}{kW}。

\textbf{8.} 压缩机：$c_p\ln(325/318) \approx \qty{0.022}{kJ/kg/K}$；
节流阀：$s_4 - s_3 = 1.204 - 1.19 = \qty{0.014}{kJ/kg/K}$——压缩机的更大。

\textbf{9.} $313/55 = 5.7$；压力比几乎翻倍（每千克做功更多），
而蒸气密度减半（每冲程质量更少）：
制冷量恰在房屋需要更多时下降。

\textbf{10.} 潮湿空气中线圈温度低于 \qty{0}{\celsius}；霜
层起隔热作用；除霜时循环反向运行，每小时数分钟，
耗费约百分之几。

\textbf{11.} $308/35 = 8.8$，对比 $328/55 = 6.0$：冷凝温度越低，
性能越好——地板下供暖。

\textbf{12.} \qty{1}{kWh} 燃料 $\to$ \qty{0.4}{kWh} 电能 $\to$
\qty{1.6}{kWh} 热量，锅炉则为 \qty{0.9}{kWh}。

\textbf{13.} 泵、锅炉（等压加热、汽化、过热），
涡轮机、冷凝器。

\textbf{14.} \qty{6}{kJ/kg}；$h_2 = 144$。

\textbf{15.} $3278 = 1069 + 1571 + 638$（\unit{kJ/kg}）。

\textbf{16.} $x = 0.81$，$h_{4s} = 2098$，$w_{\text{s}} = \qty{1324}{kJ/kg}$。

\textbf{17.} \qty{1125}{kJ/kg}，$h_4 = 2297$，$x_4 = 0.89$；较湿的蒸汽会侵蚀
叶片并降低效率。

\textbf{18.} $\eta = 1119/3278 = 0.34$；卡诺效率 $0.60$；$\bar T = 3278/6.40 =
\qty{512}{K}$，得 $0.40$。

\textbf{19.} \qty{536}{kg/s}；\qty{1160}{MW}；\qty{28}{m^3/s} 的河水，
或 \qty{480}{kg/s} 被蒸发。

\textbf{20.} \qty{0.65}{kJ/kg/K}，\qty{350}{kW/K}；\qty{200}{kJ/kg}，占
涡轮机功的 $18\%$。

\textbf{21.} 第二次吸热处于高温（平均温度
$\bar T$ 更高），第二次膨胀始于过热状态，终态更干燥。

\textbf{22.} $1757\,\unit{MW} \times (1/512 - 1/1500) = \qty{2.3}{MW/K}$——约为
涡轮机的六倍。

\textbf{23.} $1160\,\unit{MW} \times (1/293 - 1/306) = \qty{0.17}{MW/K}$。

\textbf{24.} 炉膛 $\gg$ 涡轮机 $>$ 冷凝器：因此需采用更高的蒸汽
温度与压强，并采用再热与回热；其次是更好的
叶片；冷凝器则最后考虑。

\textbf{25.} 画出\omterm{def:b2:open-systems:cv}{控制体}；流入质量等于流出质量；$\Delta h +
\Delta c^2/2 + g\Delta z = w_{\text{u}} + q$，其中 $h$ 由 $c_p$、
表格或 $v\Delta P$ 确定；$\Delta s = \sum q/T + s_{\text{c}}$，借助等熵参考态与
效率；组合各部件并与卡诺循环比较。
\end{solution}
